% !TEX root = chapter-standalone.tex


\section{Dynamical systems}
\label{sec:dynamical-systems}

Dynamical systems are ubiquitous in engineering, and there are various ways to formulate this concept mathematically.
At a very basic level, there are usually the following ingredients:
\begin{enumerate}
    \item a \emph{state space} which comprises the possible states of the system in question
    \item a \emph{dynamics rule} which governs how states of the system may change
\end{enumerate}
In this section, there is a model of \emph{time} involved (e.g. either in discrete time-steps, or as continuous time) for describing how states change over time.
(If, instead, changes of state are triggered by \emph{events}, and only the order of these events matters, we speak of \emph{machines}; see \cref{sec:machines}.)

A common way to model dynamics is by an action of a monoid $\mathbf{T} = \tup{\stylesets{T}, \mthen_{\mathbf{T}}, \idmon_{\mathbf{T}}}$, which models time, on a set $\stylesets{X}$, which models the state space.
In uncurried form (\cref{def:uncurried-action}), this is a function
\begin{equation}
    \dyn \colon  \stylesets{T} \cartprod \stylesets{X} \sto \stylesets{X}
\end{equation}
such that $\dyn(s \mthen_{\mathbf{T}} t, x) = \dyn(t, \dyn(s, x))$ and $\dyn(\idmon_{\mathbf{T}}, x) = x$ for all $s, t \setin \stylesets{T}$ and $x \setin \stylesets{X}$: letting time $s$ pass and then time $t$ is the same as letting time $s \mthen_{\mathbf{T}} t$ pass.


\subsection{Discrete-time systems}

Consider for example the monoid $\mathbf{T} = \tup{\natnumbers, +, 0}$ as a discrete model of time.
Observe that $\mathbf{T}$ is generated by the number $1 \setin \natnumbers$: every natural number $n$ can be written as a $n$-fold sum of $1$ with itself.
If we model a dynamical system by an action of $\mathbf{T}$ on a state space $\stylesets{X}$, in uncurried form
\begin{equation}\label{eq:discrete-time-dynamics}
    \dyn \colon  \natnumbers \cartprod \stylesets{X}   \sto \stylesets{X},
\end{equation}
then in particular this dynamics function is compatible with the summing of natural numbers (which is the composition operation in the monoid $\mathbf{T}$):
\begin{equation}
    \dyn(m + n, x) = \dyn(n, \dyn(m, x)).
\end{equation}
Moreover, since $0$ is the neutral element of $\mathbf{T}$, we have $\dyn(0, x) = x$ for all $x \setin \stylesets{X}$: letting no time pass leaves the state unchanged.


\begin{example}[A clock]
    \label{exa:clock-dynamical-system}
    Let $\stylesets{X} = \makeset{0, 1, \dots, 11}$ be the set of hours shown on a clock face, and let
    \begin{equation}
        \dyn \colon \natnumbers \cartprod \stylesets{X} \sto \stylesets{X}, \qquad \dyn(n, x) \definedas (x + n) \bmod 12,
    \end{equation}
    the hour shown $n$ hours after the clock showed $x$ (here $k \bmod 12$ denotes the remainder of $k$ after division by $12$).
    This is an action of $\tup{\natnumbers, +, 0}$: we have $\dyn(0, x) = x$, and
    \begin{equation}
        \begin{aligned}
            \dyn(n, \dyn(m, x)) & = \bigl(((x + m) \bmod 12) + n\bigr) \bmod 12 \\
                                & = (x + m + n) \bmod 12 = \dyn(m + n, x),
        \end{aligned}
    \end{equation}
    which expresses the familiar fact that waiting $m$ hours and then $n$ more hours is the same as waiting $m + n$ hours.
    The state space is finite, and the dynamics is periodic: $\dyn(12, -) = \catidat{\stylesets{X}}$.
    The days of the week, with $\dyn(n, x) = (x + n) \bmod 7$, form a system of the same kind.
\end{example}

\begin{example}[Compound interest]
    \label{exa:compound-interest}
    Consider a savings account with an interest rate $r \geq 0$ per year, and no deposits or withdrawals.
    If the balance is $x \setin \reals$ at some point, then the balance after $n$ years is
    \begin{equation}
        \dyn(n, x) \definedas (1 + r)^n \, x.
    \end{equation}
    This is an action of $\tup{\natnumbers, +, 0}$ on $\reals$, since $\dyn(0, x) = (1 + r)^0 x = x$ and
    \begin{equation}
        \dyn(n, \dyn(m, x)) = (1 + r)^n (1 + r)^m \, x = (1 + r)^{m + n} \, x = \dyn(m + n, x).
    \end{equation}
    Here the state space is a continuum, and the dynamics is linear in the state.
    We will return to this system in \cref{exa:savings-account}, where we also allow deposits and withdrawals.
\end{example}


\begin{counterexample}[Quadratic growth]
    \label{cexa:quadratic-growth}
    Suppose that some quantity increases by $n^2$ in $n$ time steps, as described by the function
    \begin{equation}
        \dyn \colon \natnumbers \cartprod \reals \sto \reals, \qquad \dyn(n, x) \definedas x + n^2.
    \end{equation}
    The condition $\dyn(0, x) = x$ holds, but compatibility with composition fails: letting one time step pass and then another gives $\dyn(1, \dyn(1, x)) = x + 2$, whereas letting two time steps pass at once gives $\dyn(2, x) = x + 4$.
    So $\dyn$ is not an action of $\tup{\natnumbers, +, 0}$.
    The problem is that the increase during the next time step, $(n + 1)^2 - n^2 = 2n + 1$, depends on how much time has already passed, and this information is not part of the state $x$.
    It can be repaired by including the elapsed time in the state: on $\stylesets{X} = \makecartprod{\reals, \natnumbers}$,
    \begin{equation}
        \dyn(n, \tup{x, k}) \definedas \tup{x + (k + n)^2 - k^2, \, k + n}
    \end{equation}
    is an action of $\tup{\natnumbers, +, 0}$, as one may check by a short computation.
    In general, the state of a dynamical system must contain all the information that is needed to predict its future.
\end{counterexample}

\paragraph{The one-step function of a discrete-time system.}
Because the monoid $\mathbf{T}$ is generated by the number $1$, the function \cref{eq:discrete-time-dynamics} is completely determined by the function
\begin{equation}\label{eq:discrete-dynam-generating}
    \dyn(1, - ) \colon \stylesets{X}   \sto \stylesets{X}.
\end{equation}
We call this function the \maindef{one-step function} of the discrete-time dynamical system.
If we start ``running'' the dynamical system at some point $x_0$ at time $t = 0$, then iteratively applying the function \cref{eq:discrete-dynam-generating} will compute the whole trajectory $x_0, x_1, x_2, x_3, \dots$ of the initial point $x_0$ as time progresses.

\begin{lemma}
\label{lem:one-step-function}
For any set $\stylesets{X}$, there is a bijection from the set of (uncurried) actions
\begin{equation}
\natnumbers \cartprod \stylesets{X}  \sto \stylesets{X}
\end{equation}
of $\tup{\natnumbers, +, 0}$ to the set of functions $\mora \colon \stylesets{X} \sto \stylesets{X}$; it is given by sending any action $\dyn(- , -)$ to its one-step function $\dyn(1 , -)$.
\end{lemma}

\begin{remark}
\label{rem:one-step-function-inverse}
The inverse of the bijection of \cref{lem:one-step-function} sends a function $\mora \colon \stylesets{X} \sto \stylesets{X}$ to the action
\begin{equation}
    \natnumbers \cartprod \stylesets{X} \sto \stylesets{X}, \qquad \tup{n, x} \mapsto \mora^n(x),
\end{equation}
where $\mora^n = \mora \mthen \cdots \mthen \mora$ is the $n$-fold composite of $\mora$ with itself, and $\mora^0 \definedas \catidat{\stylesets{X}}$.
\end{remark}

\begin{example}[Plant states]
    \label{exa:plant-dynamical-system}
    Recall the set $\setA = \makeset{\sprout, \yng, \mature, \old, \dead}$ of states of a plant and the development map $\mapa \colon \setA \sto \setA$ of \cref{exa:plant-trafo-semigroup} (see \cref{fig:plants_transitions}), which describes the change of state during one season:
    \begin{equation}
        \sprout \mapsto \yng \mapsto \mature \mapsto \old \mapsto \dead, \qquad \dead \mapsto \dead.
    \end{equation}
    Taking $\mapa$ as the one-step map \cref{eq:discrete-dynam-generating}, with $\mapa^0 \definedas \catidat{\setA}$, we obtain a discrete-time dynamical system
    \begin{equation}
        \dyn \colon \natnumbers \cartprod \setA \sto \setA, \qquad \dyn(n, \statesetel) \definedas \mapa^n(\statesetel),
    \end{equation}
    which describes the state of the plant after $n$ seasons.
    It is an action of $\tup{\natnumbers, +, 0}$, since $\mapa^{m+n} = \mapa^m \mthen \mapa^n$ and $\mapa^0 = \catidat{\setA}$.
    (In \cref{exa:plant-trafo-semigroup} we only considered the powers $\mapa^n$ with $n \geq 1$, which form a semigroup; allowing $n = 0$ turns this into an action of a monoid.)

    The trajectory starting at $\sprout$ is
    \begin{equation}
        \sprout, \ \yng, \ \mature, \ \old, \ \dead, \ \dead, \ \dead, \ \dots
    \end{equation}
    Every trajectory reaches the state $\dead$ after at most four seasons and stays there: $\dead$ is a \emph{fixed point} of the dynamics, $\mapa(\dead) = \dead$.
    Accordingly, the functions $\dyn(n, -)$ are not all different: $\mapa^0, \mapa^1, \dots, \mapa^4$ are five different functions, and $\mapa^n = \mapa^4$ for all $n \geq 4$, since $\mapa^4$ sends every state to $\dead$.
    Despite its simplicity, this system illustrates the basic ingredients: a state space, a one-step map, and the trajectories obtained by iterating it.
\end{example}




Often a discrete-time dynamical system is specified by defining the one-step function in the following notation,
\begin{equation}\label{eq:discrete-recursion-dynamics}
    x_{n+1} = \mapa(x_n),
\end{equation}
where $\mapa = \dyn(1, - )$ and $x_n$ denotes $\dyn(n, x_0)$, the $n$-th element of a trajectory starting at some $x_0$.

\begin{example} The logistic equation
    \begin{equation}
        x_{n+1} = rx_n(1 - x_n)
    \end{equation}
    is of the form \cref{eq:discrete-recursion-dynamics}, with $\stylesets{X} = [0,1] \subseteq \reals$ and where $r \setin [0,4]$ is a parameter (this range ensures that the map sends $[0,1]$ to itself).
    Despite being only 1-dimensional and deterministic, this dynamical system can exhibit extremely complex behavior (more precisely: deterministic chaos), depending on which value for the parameter $r$ is chosen.
\end{example}



\paragraph{The graph of a finite discrete-time system.}
We call a dynamical system \maindef{finite} if its state space is a finite set.
For a finite discrete-time system, we can draw its one-step function $\mora = \dyn(1, -)$ as a graph: each element of the state space $\stylesets{X}$ is represented by a dot, and for each $x \setin \stylesets{X}$ we draw an arrow from the dot of $x$ to the dot of $\mora(x)$, indicating where $x$ is sent in one time step.
Since $\mora$ is a function, exactly one arrow starts at each dot; this arrow may lead back to the dot it starts from (a \emph{loop}), and a dot may be the end of any number of arrows, including none.
By \cref{lem:one-step-function}, this graph contains all the information about the system: the state $\dyn(n, x)$ is reached by starting at the dot of $x$ and following $n$ arrows, and the trajectory of $x$ is the sequence of dots visited along the way.

\begin{example}[The graph of the plant system]
    \label{exa:plant-graph}
    \begin{marginfigure}
        \centering
        \begin{tikzpicture}[y=-0.95cm,
                dot/.style={circle, fill, inner sep=1.4pt},
                arr/.style={-{Stealth}, shorten <=2pt, shorten >=2pt}]
            \node[dot, label=right:{$\sprout$}] (s) at (0,0) {};
            \node[dot, label=right:{$\yng$}]    (y) at (0,1) {};
            \node[dot, label=right:{$\mature$}] (m) at (0,2) {};
            \node[dot, label=right:{$\old$}]    (o) at (0,3) {};
            \node[dot, label=right:{$\dead$}]   (d) at (0,4) {};
            \draw[arr] (s) -- (y);
            \draw[arr] (y) -- (m);
            \draw[arr] (m) -- (o);
            \draw[arr] (o) -- (d);
            \draw[arr] (d) .. controls +(-0.9,0.9) and +(0.9,0.9) .. (d);
        \end{tikzpicture}
        \caption{The graph of the one-step function of the plant system (\cref{exa:plant-dynamical-system}).}
        \label{fig:plant-graph}
    \end{marginfigure}
    The graph of the one-step function of the plant system of \cref{exa:plant-dynamical-system} is shown in \cref{fig:plant-graph}.
    It is a single chain of arrows ending in a loop at $\dead$.
    Several features of the system can be read off directly: the loop shows that $\dead$ is a fixed point; following arrows from any dot, we reach $\dead$ after at most four steps; no arrow ends at $\sprout$, so $\sprout$ is not the successor of any state; and two arrows end at $\dead$, so the one-step function is not injective, and the past of a plant cannot always be recovered from its present.
\end{example}

\begin{example}[A traffic light]
    \label{exa:traffic-light}
    \begin{marginfigure}
        \centering
        \begin{tikzpicture}[
                dot/.style={circle, fill, inner sep=1.4pt},
                arr/.style={-{Stealth}, shorten <=2pt, shorten >=2pt}]
            \node[dot, label=above:{$\text{start}$}] (st) at (-1.6,0) {};
            \node[dot, label=above:{$\text{red}$}]    (r)  at (0,0) {};
            \node[dot, label=above:{$\text{green}$}]  (g)  at (1.6,0) {};
            \node[dot, label=below:{$\text{yellow}$}] (yl) at (0.8,-1.3) {};
            \draw[arr] (st) -- (r);
            \draw[arr] (r) -- (g);
            \draw[arr] (g) -- (yl);
            \draw[arr] (yl) -- (r);
        \end{tikzpicture}
        \caption{The graph of the one-step function of the traffic light (\cref{exa:traffic-light}).}
        \label{fig:traffic-light-graph}
    \end{marginfigure}
    A traffic light is switched on in a state $\text{start}$, and then goes through the phases red, green and yellow, after which it starts again with red.
    With one time step per phase, this is the finite discrete-time system with state space $\stylesets{X} = \makeset{\text{start}, \text{red}, \text{green}, \text{yellow}}$ and one-step function
    \begin{equation}
        \text{start} \mapsto \text{red}, \quad \text{red} \mapsto \text{green}, \quad \text{green} \mapsto \text{yellow}, \quad \text{yellow} \mapsto \text{red}.
    \end{equation}
    Its graph (\cref{fig:traffic-light-graph}) has a different shape from that of the plant system: a ``tail'' consisting of the single arrow from $\text{start}$ to $\text{red}$, followed by a cycle $\text{red} \to \text{green} \to \text{yellow} \to \text{red}$ of length $3$.
    Accordingly, the trajectory starting at $\text{start}$,
    \begin{equation}
        \text{start}, \ \text{red}, \ \text{green}, \ \text{yellow}, \ \text{red}, \ \text{green}, \ \text{yellow}, \ \text{red}, \ \dots,
    \end{equation}
    never comes back to $\text{start}$, but from time $1$ on it is periodic, with period $3$: $\dyn(n + 3, x) = \dyn(n, x)$ for all $n \geq 1$ and all $x \setin \stylesets{X}$.
\end{example}

For finite discrete-time dynamical systems, every trajectory must eventually visit some state for a second time, and from then on the trajectory begins to repeat itself. Thus, every trajectory consists of a (possibly empty) tail, followed by a cycle. The clock of \cref{exa:clock-dynamical-system}, for instance, has no tails at all: its graph is a single cycle of length $12$. A loop, as at $\dead$ in the plant system, is a cycle of length $1$.


\paragraph{Integer time.}
Instead of the natural numbers, we can also model discrete time by the integers, that is, by the monoid $\tup{\wnumbers, +, 0}$, which is in fact a group (the inverse of $n$ is $-n$).
A dynamical system with integer time is an action of $\tup{\wnumbers, +, 0}$ on a set $\stylesets{X}$, in uncurried form
\begin{equation}
    \label{eq:integer-time-dynamics}
    \dyn \colon \wnumbers \cartprod \stylesets{X} \sto \stylesets{X}.
\end{equation}
Its trajectories $x_n = \dyn(n, x_0)$ extend not only into the future but also into the past: $\dots, x_{-2}, x_{-1}, x_0, x_1, x_2, \dots$.

As before, such a system is determined by the one-step map $\mapa = \dyn(1, -)$, since every integer is a finite sum of copies of $1$ or of $-1$, and the map $\dyn(-1, -)$ is in turn determined by $\mapa$: from $1 + (-1) = 0 = (-1) + 1$ we get
\begin{equation}
    \mapa \mthen \dyn(-1, -) = \dyn(0, -) = \catidat{\stylesets{X}} = \dyn(-1, -) \mthen \mapa.
\end{equation}
So $\mapa$ must be invertible, with inverse $\dyn(-1, -)$.
Conversely, every invertible function $\mapa \colon \stylesets{X} \sto \stylesets{X}$ determines an action of $\tup{\wnumbers, +, 0}$, by sending $n \geq 0$ to $\mapa^n$ and $n < 0$ to $(\mapa^{-1})^{-n}$.
In other words, a dynamical system with integer time is the same as a discrete-time system whose dynamics is \emph{reversible}: from the current state we can compute not only the next state, $x_{n+1} = \mapa(x_n)$, but also the previous one, $x_{n-1} = \mapa^{-1}(x_n)$.
Since every $\dyn(n, -)$ is then invertible, such a system is also the same as a \SY{group action} of $\tup{\wnumbers, +, 0}$ on $\stylesets{X}$, that is, a group morphism $\wnumbers \fto \Autof{\stylesets{X}}$ (\cref{def:group-cov-action}).

The two models of discrete time are related as follows.
Every system with integer time gives a system with time $\tup{\natnumbers, +, 0}$, by forgetting about negative times (formally: by restricting along the inclusion $\natnumbers \sto \wnumbers$, which is a monoid morphism; see \cref{lem:restriction-of-action} below).
Conversely, a system with time $\tup{\natnumbers, +, 0}$ and one-step map $\mapa$ extends to a system with integer time exactly if $\mapa$ is invertible, and in that case the extension is unique.
For example, the linear system $x_{n+1} = \mat{A} x_n$ on $\reals^d$ extends to integer time exactly if the $\matdim{d}{d}$ matrix $\mat{A}$ is invertible.
The logistic map, on the other hand, is not injective for $r > 0$ (it sends $x$ and $1 - x$ to the same value), so it defines a system with time $\tup{\natnumbers, +, 0}$ only: different pasts can lead to the same present, and the past cannot be recovered from the present.
The same holds for the plant model of \cref{exa:plant-dynamical-system}, where both $\old$ and $\dead$ are sent to $\dead$.
(The flows of the next subsection are the continuous-time counterpart of integer time: they are actions of $\tup{\reals, +, 0}$, defined for negative as well as positive times.)

\subsection{Continuous-time systems}

Consider the initial value problem
\begin{equation}
    \begin{cases}
        \dot{\stylemaps{x}}(t) = \mapa(\stylemaps{x}(t)) \\
        \stylemaps{x}(0) = x_0
    \end{cases}
\end{equation}
where $\mapa \colon \reals^n \sto \reals^n$ is a differentiable function, $\stylemaps{x} \colon \reals \sto \reals^n$ is a continuously differentiable function, and $x_0 \setin \reals^n$.
Under suitable assumptions on $\mapa$, for each choice of $x_0$, there exists a unique solution $\stylemaps{s}_{x_0} \colon \reals \sto \reals^n$ of this initial value problem.
Using all of these solutions, it is possible to define the following map, which is called a \emph{flow}:
\begin{equation}
    \defmapperiod{
        \varphi
    }{
        \reals \cartprod \reals^n
    }{
        \to
    }{
        \reals^n
    }{
        \tup{t,x_0}
    }{
        \stylemaps{s}_{x_0}(t)
    }
\end{equation}

In this case, the flow map $\varphi$ will necessarily satisfy the following conditions
\begin{enumerate}
    \item $\varphi(t + s, x) = \varphi(s, \varphi(t,x)) \quad \forall t, s \setin \reals, x \setin \reals^n$,
    \item $\varphi(0, x) = x \quad \forall x \setin \reals^n$,
\end{enumerate}
which means that it is an (uncurried) action of the monoid $\tup{\reals, +, 0}$ on $\reals^n$.

\todotext{JL: explain in a couple more words the intuition/idea, and perhaps include a graphic here}

\begin{definition}
    \label{def:ct-dynsyst}
    A continuous-time dynamical system on $\reals^n$ is:

    \constit
    \begin{enumerate}
        \item A continuously differentiable function
              \begin{equation}
                  \dyn \colon \reals \cartprod \reals^n \to \reals^n.
              \end{equation}
    \end{enumerate}

    \condit
    \begin{enumerate}
        \item $\dyn$ is an action of the monoid $\tup{\reals, +, 0}$.
    \end{enumerate}
\end{definition}

\subsection{Sampling a continuous-time system}
\label{sec:sampling}

In practice, a continuous-time system is often only observed, or controlled, at discrete instants of time -- for example, when a digital sensor reads the state every $h$ seconds.
The discrete-time system obtained in this way arises from the continuous-time one by a general construction, which works for actions of arbitrary monoids: an action can be \emph{restricted along a monoid morphism}.

\begin{lemma}
    \label{lem:restriction-of-action}
    Let $\monoidA$ and $\monoidB$ be \SY{monoids}, let $\act \colon \monoidA \fto \Endof\stateset$ be an action of $\monoidA$ on a set $\stateset$, and let $\stylefunctors{\sigma} \colon \monoidB \fto \monoidA$ be a \SY{monoid morphism}.
    Then the composite function
    \begin{equation}
        \stylefunctors{\sigma} \mthen \act \colon \monoidB \fto \Endof\stateset
    \end{equation}
    is an action of $\monoidB$ on $\stateset$.
    In uncurried form, it is given by $\tup{\monelA, \statesetel} \mapsto \act(\stylefunctors{\sigma}(\monelA), \statesetel)$.
\end{lemma}

\begin{proof}
    We must check that $\stylefunctors{\sigma} \mthen \act$ is a \SY{monoid morphism}.
    For all $\monelA, \monelB \setin \monoidBset$ we have
    \begin{equation}
        \act(\stylefunctors{\sigma}(\monelA \mtimesof\monoidB \monelB))
        = \act(\stylefunctors{\sigma}(\monelA) \mtimesof\monoidA \stylefunctors{\sigma}(\monelB))
        = \act(\stylefunctors{\sigma}(\monelA)) \mtimesof{\Endof\stateset} \act(\stylefunctors{\sigma}(\monelB)),
    \end{equation}
    using first that $\stylefunctors{\sigma}$ and then that $\act$ is compatible with composition; similarly, $\act(\stylefunctors{\sigma}(\idmon_\monoidB)) = \act(\idmon_\monoidA) = \catidat\stateset$.
\end{proof}

Now fix a \emph{sampling period} $h \setin \posReals$.
The function
\begin{equation}
    \defmapperiod{\stylefunctors{\sigma}_h}{\natnumbers}{\to}{\reals}{k}{k h}
\end{equation}
is a \SY{monoid morphism} $\tup{\natnumbers, +, 0} \fto \tup{\reals, +, 0}$, since $(k + k') h = k h + k' h$ and $0 h = 0$.
Given a continuous-time dynamical system $\dyn \colon \reals \cartprod \reals^n \sto \reals^n$ (\cref{def:ct-dynsyst}), \cref{lem:restriction-of-action} therefore yields an action of $\tup{\natnumbers, +, 0}$ on $\reals^n$, that is, a discrete-time dynamical system:
\begin{equation}
    \dyn_h \colon \natnumbers \cartprod \reals^n \sto \reals^n, \qquad \dyn_h(k, x) \definedas \dyn(k h, x).
\end{equation}
This is the \emph{sampled system}: $\dyn_h(k, x_0)$ is the state at time $k h$ of the trajectory starting at $x_0$.
As for any discrete-time system, it is determined by its one-step map $\dyn_h(1, -) = \dyn(h, -)$, and in recursion notation \cref{eq:discrete-recursion-dynamics} it reads
\begin{equation}
    x_{k+1} = \dyn(h, x_k).
\end{equation}

\begin{example}
    \label{exa:sampled-exponential}
    For the scalar differential equation $\dot{\stylemaps{x}}(t) = \lambda \stylemaps{x}(t)$, with a constant $\lambda \setin \reals$, the flow is $\dyn(t, x) = e^{\lambda t} x$.
    The sampled system with period $h$ is $x_{k+1} = e^{\lambda h} x_k$.
    More generally, for a linear system $\dot{\stylemaps{x}}(t) = \mat{A}\stylemaps{x}(t)$ on $\reals^n$, the flow is $\dyn(t, x) = e^{\mat{A}t} x$ (with the matrix exponential), and the sampled system is $x_{k+1} = e^{\mat{A}h} x_k$.
    This is the starting point for turning continuous-time LTI systems (\cref{def:LTI_syst_first_time}) into discrete-time ones (\cref{def:LTI-discrete}).
\end{example}

Sampling loses information: different continuous-time systems can have the same sampled system.
For example, the flows on $\reals^2$ which rotate every point around the origin with angular velocity $\omega$, respectively $\omega + 2\pi/h$, have the same sampled system with period $h$, since after each sampling period the two rotations differ by a full turn.
In signal processing and control, this phenomenon is known as \emph{aliasing}.

\subsection{Linear time-invariant systems}

An important class of dynamical systems in engineering are \emph{linear time-invariant (LTI) systems}.
In contrast to the systems considered so far, they have \emph{inputs} and \emph{outputs}: the state changes over time according to its current value and the current input, and the output is computed from the state and the input.
We first consider the discrete-time version, and then the continuous-time version.

\begin{definition}[Discrete-time LTI system]
    \label{def:LTI-discrete}
    A \emph{discrete-time linear time-invariant (LTI) system}, in a state-space representation, is:

    \constit
    \begin{enumerate}
        \item Real \SY{vector spaces} $\prin=\reals^l$ (input space), $\prout=\reals^m$ (output space), and $\prst=\reals^n$ (state space);
        \item Real matrices $\mat{A}$, $\mat{B}$, $\mat{C}$, $\mat{D}$ of dimensions $\matdim{n}{n}$, $\matdim{n}{l}$, $\matdim{m}{n}$, and $\matdim{m}{l}$, respectively, specifying a system of equations of the form
              \begin{equation}
                  \label{eq:LTI-discrete-equations}
                  \begin{aligned}
                      \mat{x}_{k+1} & = \mat{A}\mat{x}_k + \mat{B}\mat{u}_k, \\
                      \mat{y}_k     & = \mat{C}\mat{x}_k + \mat{D}\mat{u}_k,
                  \end{aligned}
              \end{equation}
              where $k \setin \natnumbers$, $\mat{u}_k \setin \prin$, $\mat{y}_k \setin \prout$, $\mat{x}_k \setin \prst$;
        \item An \emph{initial state} $\mat{x}_0 \setin \prst$.
    \end{enumerate}
\end{definition}

We write such a system compactly as a tuple $\tupp{\mat{A}, \mat{B}, \mat{C}, \mat{D}, \mat{x}_0}$.
Given a sequence $\mat{u}_0, \mat{u}_1, \mat{u}_2, \dots$ of inputs, the equations \cref{eq:LTI-discrete-equations} determine, step by step, the sequence of states $\mat{x}_0, \mat{x}_1, \mat{x}_2, \dots$ and the sequence of outputs $\mat{y}_0, \mat{y}_1, \mat{y}_2, \dots$; explicitly,
\begin{equation}
    \mat{x}_k = \mat{A}^k \mat{x}_0 + \sum_{i=0}^{k-1} \mat{A}^{k-1-i} \mat{B} \mat{u}_i.
\end{equation}
The first equation in \cref{eq:LTI-discrete-equations} specifies dynamics $\prdyn \colon \makecartprod{\prin, \prst} \sto \prst$, $\tup{u, x} \mapsto \mat{A}x + \mat{B}u$, which computes the \emph{next} state, and the second one specifies a read-out $\prreadout \colon \makecartprod{\prin, \prst} \sto \prout$, $\tup{u, x} \mapsto \mat{C}x + \mat{D}u$.

If the input is zero at all times, the system reduces to a discrete-time dynamical system as in \cref{sec:dynamical-systems}: the action of $\tup{\natnumbers, +, 0}$ on $\prst = \reals^n$ is given by $\dyn(k, x) = \mat{A}^k x$, with $\mat{A}^0 = \idmat$.
With inputs, the system has the same structure as a \emph{Mealy machine} (\cref{exa:discrete-lti-mealy}); the difference lies in reading the index $k$ as time.

\begin{example}[Savings account]
    \label{exa:savings-account}
    A savings account with interest rate $r$ per year, into which an amount $\mat{u}_k$ is deposited at the end of year $k$ (or withdrawn, if $\mat{u}_k$ is negative), can be modelled by the discrete-time LTI system with $l = m = n = 1$ and
    \begin{equation}
        \mat{x}_{k+1} = (1 + r) \mat{x}_k + \mat{u}_k, \qquad \mat{y}_k = \mat{x}_k,
    \end{equation}
    that is, $\mat{A} = 1 + r$, $\mat{B} = 1$, $\mat{C} = 1$, $\mat{D} = 0$, where the state $\mat{x}_k$ and the output $\mat{y}_k$ are the balance at the beginning of year $k$.
\end{example}

\begin{definition}[Continuous-time LTI system]
    \label{def:LTI_syst_first_time}
    A \emph{continuous-time linear time-invariant (LTI) system}, in a state-space representation, is:

    \constit
    \begin{enumerate}
        \item Real \SY{vector spaces} $\prin=\reals^l$ (input space), $\prout=\reals^m$ (output space), and $\prst=\reals^n$ (state space);
        \item Real matrices $\mat{A}$, $\mat{B}$, $\mat{C}$, $\mat{D}$ of appropriate dimension, specifying a system of equations of the form
              \begin{align}
                  \dot{\mat{x}}(t) & = \mat{A}\mat{x}(t) + \mat{B}\mat{u}(t) \label{eq:LTI-system-dyn-equation} \\
                  \mat{y}(t)       & = \mat{C}\mat{x}(t) + \mat{D}\mat{u}(t) \label{eq:LTI-system-ro-equation},
              \end{align}
              where $t \setin \nonNegReals$, $\mat{u}(t) \setin \prin$, $\mat{y}(t) \setin \prout$, $\mat{x}(t) \setin \prst$;
        \item An \emph{initial state} $\prstart\setin \prst$.
    \end{enumerate}
\end{definition}

We think of \cref{eq:LTI-system-dyn-equation} as defining dynamics
\begin{equation}
    \defmapcomma{
        \prdyn
    }{ \prin \cartprod  \prst
    }{ \sto
    }{  \prst
    }{ \tup{u, x}
    }{ \mat{A}x + \mat{B}u
    }
\end{equation}
and \cref{eq:LTI-system-ro-equation} as defining a read-out function
\begin{equation}
    \defmapperiod{
        \prreadout
    }{ \prin \cartprod  \prst
    }{ \sto
    }{  \prout
    }{ \tup{u, x}
    }{ \mat{C}x + \mat{D}u
    }
\end{equation}

In contrast to the discrete-time case (\cref{def:LTI-discrete}), here the inputs $\mat{u}(t)$ are functions of continuous time (``signals''), and the dynamics $\tup{u, x} \mapsto \mat{A}x + \mat{B}u$ specifies the \emph{rate of change} of the state, not the next state.
